\documentclass[12pt]{article}

% ---------------------------------------------------------
% PACKAGES
% ---------------------------------------------------------
\usepackage[a4paper,margin=0.8in]{geometry}
\usepackage{hyperref}
\usepackage{titlesec}
\usepackage{setspace}
\usepackage{enumitem}
\usepackage[nopatch=footnote]{microtype}
\usepackage{fancyhdr}
\usepackage{tabularx}
\usepackage{newtxtext,newtxmath}


% ---------------------------------------------------------
% DOCUMENT SETTINGS
% ---------------------------------------------------------
\setlength{\headheight}{14.49998pt}
\hypersetup{
    colorlinks=true,
    linkcolor=black,
    urlcolor=blue,
    citecolor=black,
    pdfauthor={Leonard Okyere Afeke},
    pdftitle={CVE-2025-15467 Cyber Threat Intelligence Report}
}

\setstretch{1.2}

\titleformat{\section}{\large\bfseries}{\thesection}{1em}{}
\titleformat{\subsection}{\normalsize\bfseries}{\thesubsection}{1em}{}

\pagestyle{fancy}
\fancyhf{}
\rhead{CVE-2025-15467 CTI Report}
\lhead{Leonard Okyere Afeke}
\cfoot{\thepage}

% ---------------------------------------------------------
% TITLE PAGE
% ---------------------------------------------------------
\begin{document}

\begin{titlepage}
    \centering
    \vspace*{3cm}
    {\Huge \textbf{CVE-2025-15467}}\par
    \vspace{0.5cm}
    {\Large Cyber Threat Intelligence Case Study}\par
    \vspace{1.5cm}

    {\large \textbf{Author:} Leonard Okyere Afeke}\par
    \vspace{0.3cm}
    {\large \textbf{Date:} \today}\par
    \vspace{2cm}

    \vfill
    {\small This report is based on publicly available information (OSINT) 
    and controlled internal telemetry collected in a lab environment. 
    No scanning, exploitation, or interaction with live systems occurred.}
\end{titlepage}

% ---------------------------------------------------------
% EXECUTIVE SUMMARY
% ---------------------------------------------------------
\section*{Executive Summary}

CVE-2025-15467 is a stack buffer overflow in OpenSSL CMS 
\texttt{AuthEnvelopedData} AEAD parsing. The vulnerability affects all 
OpenSSL 3.x branches prior to their patched releases and may enable a 
crash or remote code execution when processing a specially crafted CMS 
object.

Testing on OpenSSL 3.0.2—part of the vulnerable 3.0.x branch—showed that 
malformed CMS samples did not trigger the overflow. OpenSSL 3.x enforces 
strict ASN.1 validation, causing malformed inputs to fail early and 
preventing the vulnerable IV-copying code path from being reached.

This report applies the Cyber Threat Intelligence (CTI) lifecycle to 
analyze OSINT, internal telemetry, exploitability, and defensive 
implications.

% ---------------------------------------------------------
% KEY JUDGMENTS
% ---------------------------------------------------------
\section{Key Judgments}

\begin{itemize}
    \item OpenSSL 3.0.2 is part of the vulnerable 3.0.x branch; the 
    vulnerability is remediated only in version 3.0.19.
    \item Malformed CMS samples did not reach the vulnerable code path 
    due to strict ASN.1 validation in OpenSSL 3.x.
    \item Exploitation requires a valid AEAD CMS object with an 
    oversized IV; malformed CMS alone is insufficient.
    \item Telemetry collection confirmed full visibility into OpenSSL 
    executions, arguments, and file interactions.
    \item Behavior observed on 3.0.2 generalizes to all affected 3.x 
    branches.
\end{itemize}

% ---------------------------------------------------------
% COLLECTION
% ---------------------------------------------------------
\section{Collection}

\subsection*{OSINT Sources}
\begin{itemize}
    \item MITRE and NVD CVE entries
    \item OpenSSL security advisory
    \item Linux distribution advisories (Ubuntu, Debian, Fedora, SUSE)
    \item oss-security mailing list
    \item Patch commits across OpenSSL branches
\end{itemize}

\subsection*{Internal Telemetry}
Telemetry was collected on an Ubuntu VM running OpenSSL 3.0.2 using 
Elastic Defend. Data sources included:
\begin{itemize}
    \item process execution events
    \item file creation and access events
    \item ASN.1 parsing error output
    \item full process chains in Kibana Discover
\end{itemize}

% ---------------------------------------------------------
% PROCESSING
% ---------------------------------------------------------
\section{Processing}

OSINT and telemetry were normalized into structured fields.

\subsection*{OSINT Processing}
\begin{itemize}
    \item Vulnerability Type: Stack buffer overflow
    \item Component: CMS \texttt{AuthEnvelopedData} AEAD parsing
    \item Root Cause: Oversized ASN.1 IV copied into fixed-size buffer
    \item Affected Branches: 3.0.x, 3.3.x, 3.4.x, 3.5.x, 3.6.x
    \item Patched Versions: 3.0.19, 3.3.6, 3.4.4, 3.5.5, 3.6.1
\end{itemize}

\subsection*{Telemetry Processing}
\begin{itemize}
    \item All malformed CMS samples produced ASN.1 parsing errors.
    \item No crash or memory corruption occurred.
    \item No Elastic Defend rules triggered.
    \item OpenSSL executions were fully observable.
\end{itemize}

% ---------------------------------------------------------
% ANALYSIS
% ---------------------------------------------------------
\section{Analysis}

\subsection*{Exploitability}
OSINT indicates that exploitation requires:
\begin{itemize}
    \item a valid \texttt{AuthEnvelopedData} CMS structure
    \item AEAD cipher parameters (e.g., AES-GCM)
    \item an oversized IV embedded in ASN.1 parameters
\end{itemize}

Malformed CMS samples did not meet these conditions and were rejected 
before reaching the vulnerable code path.

\subsection*{Environmental Relevance}
OpenSSL 3.0.2 is widely deployed on Ubuntu systems. Although vulnerable, 
practical exploitation requires a highly specific CMS structure.

\subsection*{Adversary Behavior}
An adversary would need to craft a valid AEAD CMS object and deliver it 
to an application that processes untrusted CMS data. This requires 
moderate cryptographic understanding.

% ---------------------------------------------------------
% REPORTING
% ---------------------------------------------------------
\section{Reporting}

\subsection*{Impact Assessment}

\textbf{Technical Impact:}
\begin{itemize}
    \item Potential DoS or RCE if a crafted AEAD CMS object is processed.
    \item No crash occurred in testing due to early ASN.1 rejection.
\end{itemize}

\textbf{Operational Impact:}
\begin{itemize}
    \item Applications parsing untrusted CMS data are at risk.
    \item Systems using OpenSSL 1.1.1 or 1.0.2 are not affected.
\end{itemize}

\subsection*{Detection Opportunities}

\textbf{Process-Level:}
\begin{itemize}
    \item \texttt{openssl} CMS operations
    \item AEAD-related arguments
    \item abnormal exit codes
\end{itemize}

\textbf{File-Level:}
\begin{itemize}
    \item unexpected \texttt{.cms} file creation
\end{itemize}

\textbf{Application-Level:}
\begin{itemize}
    \item service crashes or restarts
    \item ASN.1 parsing errors in logs
\end{itemize}

\subsection*{Mitigation}
\begin{itemize}
    \item Upgrade to patched branch releases:
    \begin{itemize}
        \item 3.0.x $\rightarrow$ 3.0.19
        \item 3.3.x $\rightarrow$ 3.3.6
        \item 3.4.x $\rightarrow$ 3.4.4
        \item 3.5.x $\rightarrow$ 3.5.5
        \item 3.6.x $\rightarrow$ 3.6.1
    \end{itemize}
    \item Restrict CMS parsing to trusted inputs.
    \item Monitor OpenSSL activity using process and file telemetry.
\end{itemize}

% ---------------------------------------------------------
% CONCLUSION
% ---------------------------------------------------------
\section{Conclusion}

OpenSSL 3.0.2 is vulnerable to CVE-2025-15467 as part of the broader 
3.0.x branch. However, malformed CMS inputs did not reach the vulnerable 
AEAD IV-copying code path due to strict ASN.1 validation in OpenSSL 3.x. 
Exploitation requires a highly specific, well-formed AEAD CMS structure.

Telemetry confirms that OpenSSL behavior is observable and provides 
actionable detection opportunities. Upgrading to patched branch releases 
remains the primary mitigation.

\end{document}